\newpage
\section{Appendix: Liturgical Resolution}
All Saints is a solemnity, white, in the U.S. national calendar on Sunday 1 November 2026, replacing the Thirty-first Sunday in Ordinary Time. The surrounding Sunday cycle is A; the All Saints readings are annual, common to A, B and C. The weekday cycle is inapplicable. The governing books are the U.S. \work{Roman Missal}, Third Edition (2011), with the amended third typical Latin edition, and the U.S. second typical \work{Lectionary for Mass}, no. 667. An anticipated celebration uses this Mass; there is no separately appointed Vigil Mass.

The identity is PC-R07-ABC. Its reusable Missal evidence belongs to the same edition's All Saints shared owner. Gloria and Creed are said. The proper Preface is retained with the Roman Canon or Eucharistic Prayer III, or with II under its stated conditions. Eucharistic Prayer IV's invariable Preface makes it inapplicable here. The Missal antiphons are treated without asserting an actual singing choice. General chant permissions and Sunday sprinkling remain open possibilities; no particular parish selection or local ritual overlay is established. There is no appointed Sequence or festival-specific Eucharistic Prayer insert. Solemn Blessing no. 18 is permitted if selected. The expansive study's resolution appendix supplies the complete branch table and authorities.

\section{Appendix: Scope and Qualifications}
This concise study derives its two interpretations from the canonical expansive All Saints study and the same owner's reviewed research. The comparisons among the Mass's texts are theological synthesis. Augustine, Aquinas, Bede and Chrysostom are cited for their comments on the passages they explain, rather than as authors of a complete interpretation of this modern Mass. The canonical instance, proper audit, research scope, interpretations, source bindings and historical appendix retain the complete evidence and limits.

Biblical wording cited here follows the historical Douay--Rheims/Challoner study translation, public domain in the United States. It is not the approved text proclaimed at the celebration; the U.S. Lectionary controls that English. NPNF, ANF and Edward Marshall's Bede translation retain their historical public-domain status. The identified English patristic deliveries were checked without a new original-language collation. Aquinas's electronic Latin Matthew deliveries name two different printed textual bases, kept distinct in References.

This is a rights-limited study. Full protected prayer and composed-antiphon bodies are withheld. The secondary 2002 Latin witness and official 2008 change list do not authenticate every word of the complete 2008 altar book. A current authoritative U.S. publisher exemplar for every oration and the blessing remains unavailable after the recorded second search. The secondary English witness has unfinished references and an older Collect conclusion. Official antiphon and Preface excerpts have narrower authority for their own texts. Commentary on their subjects supplies no substitute translation. Chronological alternatives and place qualifications follow the canonical generated record and its expansive historical appendix. The national-calendar and branch conclusions are bounded to the stated territory and date.

\section{References}
\begin{itemize}[leftmargin=1.3em,itemsep=.3em]
\item \textbf{Scripture.} \work{Douay--Rheims Bible}, Challoner text, Project Gutenberg eBook 1581, retained verse-text edition. The appointments and contextual passages are identified at use.
\item \textbf{Augustine.} \work{Our Lord's Sermon on the Mount}, I.1--4, §§1--12, trans. William Findlay; \work{Tractates on the First Epistle of John}, IV.4--7,9, trans. H. Browne; \work{Expositions on the Psalms}, Psalm 23, §§1--6. NPNF, first series, VI--VIII, retained CCEL historical English deliveries. The English psalm heading ``XXIV'' explicitly identifies Latin XXIII.
\item \textbf{Thomas Aquinas.} \work{Super Psalmos} 23, nn. 2--4, paragraphs 86989--86991, \url{https://www.corpusthomisticum.org/cps21.html}. \work{Super Matthaeum}, reportatio of Peter of Andria, cap. 5, lectio 2, paragraph 87340, \url{https://www.corpusthomisticum.org/cma0500.html}, header identifying Turin 1951; continuation, paragraphs 87564--87566, \url{https://www.corpusthomisticum.org/cma0510.html}, header identifying Leonine 1969. These are electronic Latin witnesses, without a claimed print-image collation.
\item \textbf{Bede.} \work{Explanation of the Apocalypse}, trans. Edward Marshall (Oxford and London: James Parker, 1878), I, on 7:2,4,9--14, pp. 44--53. His biblical catchwords are distinct from appointed English.
\item \textbf{John Chrysostom.} \work{Homilies on the Gospel of St. Matthew}, XV.2--9 and XXXVIII.3, trans. George Prevost, revised M. B. Riddle, NPNF, first series, X, retained CCEL historical English delivery.
\item \textbf{Received Victorinus commentary.} \work{Commentary on the Apocalypse}, seventh chapter, comments 2,9, trans. Robert Ernest Wallis, ANF VII. Its edition's elucidations identify the recension and attribution difficulties.
\item \textbf{Liturgical controls.} U.S. \work{Roman Missal}, Third Edition (2011), All Saints and Solemn Blessing no. 18; qualified secondary 2002 Latin reproduction, PDF 512--514,377--378, with the official 2008 \work{Variationes}, \work{Notitiae} 503--504, pp. 368--387. ICEL \work{Antiphonary}, All Saints, p. 134; \href{https://www.icelweb.org/musicfolder/73JerusalemD.pdf}{ICEL All Saints Preface}, pp. 1--2. \work{Ordo lectionum Missae} (1981), no. 667, p. 313. USCCB \work{Liturgical Calendar 2026}, pp. 5,42; \href{https://bible.usccb.org/bible/readings/110126.cfm}{daily readings for 1 November 2026}; U.S. GIRM, nos. 48,51,61--62,74,87,365. Full locators and source qualifications are in the canonical study and audits.
\item \textbf{Historical orientation.} USCCB, NABRE, ``The Book of Psalms,'' introduction; E. Jacquier, ``Gospel of St. Matthew,'' \work{Catholic Encyclopedia} X (1911); Walter Drum, ``Epistles of Saint John,'' VIII (1910); A. Durand, ``The New Testament,'' XIV (1912); F. Van den Biesen, ``Apocalypse,'' I (1907). Exact date/place discussions are explained in the expansive historical appendix named on page 2.
\end{itemize}
